\documentclass[11pt,a4paper]{article}

% ── Packages ──
\usepackage[utf8]{inputenc}
\usepackage[T1]{fontenc}
\usepackage{geometry}
\usepackage{setspace}
\usepackage{booktabs}
\usepackage{longtable}
\usepackage{tabularx}
\usepackage{xcolor}
\usepackage{listings}
\usepackage{hyperref}
\usepackage{fancyhdr}
\usepackage{titlesec}
\usepackage{enumitem}
\usepackage{microtype}

% ── Page Geometry ──
\geometry{margin=2.5cm, bottom=3cm}
\setlength{\parindent}{0pt}
\setlength{\parskip}{6pt}
\onehalfspacing

% ── Colors ──
\definecolor{primary}{HTML}{6366f1}
\definecolor{text}{HTML}{1e293b}
\definecolor{muted}{HTML}{64748b}
\definecolor{codebg}{HTML}{f8f9fb}
\definecolor{border}{HTML}{e4e7ee}

% ── Hyperlinks ──
\hypersetup{
    colorlinks=true,
    linkcolor=primary,
    urlcolor=primary,
    citecolor=primary
}

% ── Header / Footer ──
\pagestyle{fancy}
\fancyhf{}
\fancyhead[L]{\small\textcolor{muted}{Depression Risk Screening}}
\fancyhead[R]{\small\textcolor{muted}{Project Report}}
\fancyfoot[C]{\thepage}
\renewcommand{\headrulewidth}{0.4pt}
\renewcommand{\footrulewidth}{0pt}

% ── Section Formatting ──
\titleformat{\section}{\Large\bfseries\color{primary}}{}{0em}{}[\vspace{-4pt}\rule{\textwidth}{1pt}]
\titleformat{\subsection}{\large\bfseries\color{text}}{}{0em}{}
\titleformat{\subsubsection}{\normalsize\bfseries\color{text}}{}{0em}{}

% ── Code Listings ──
\lstdefinestyle{py}{
    language=Python,
    backgroundcolor=\color{codebg},
    basicstyle=\small\ttfamily,
    breaklines=true,
    frame=single,
    rulecolor=\color{border},
    numbers=none,
    showstringspaces=false,
    keywordstyle=\color{primary},
    commentstyle=\color{muted}
}
\lstdefinestyle{bash}{
    backgroundcolor=\color{codebg},
    basicstyle=\small\ttfamily,
    breaklines=true,
    frame=single,
    rulecolor=\color{border}
}

% ── Custom Commands ──
\newcommand{\field}[2]{\textbf{#1} & #2 \\ \hline}
\newcommand{\riskrow}[2]{#1 & #2 \\ \hline}

% ── Document ──
\begin{document}

% ══════════════════════════════════════════════════════════════════
% TITLE PAGE
% ══════════════════════════════════════════════════════════════════
\begin{titlepage}
\centering
\vspace*{4cm}
{\Huge\bfseries\color{primary} Depression Risk Screening}\\[12pt]
{\Large ML-Powered Web Application}\\[24pt]
\rule{0.6\textwidth}{1pt}\\[24pt]
{\large Project Report}\\[12pt]
{\normalsize\textcolor{muted}{May 2026}}\\[4cm]

\vfill
{\normalsize\textcolor{muted}{Student Depression Dataset \;$\cdot$\; Ensemble ML \;$\cdot$\; Flask Web App}}
\end{titlepage}

% ══════════════════════════════════════════════════════════════════
% TABLE OF CONTENTS
% ══════════════════════════════════════════════════════════════════
\newpage
\tableofcontents
\newpage

% ══════════════════════════════════════════════════════════════════
% 1. PROJECT OVERVIEW
% ══════════════════════════════════════════════════════════════════
\section{Project Overview}

A production-grade clinical depression risk screening tool built with an ensemble machine learning pipeline served through a Flask web interface. The application accepts 16 raw input features across demographic, academic, and lifestyle categories, engineers 16 additional derived features, and returns a depression risk probability via a soft-voting ensemble of four classifiers.

Key highlights:
\begin{itemize}[nosep]
    \item \textbf{Dataset}: ~27,900 clinical records (Kaggle Student Depression Dataset)
    \item \textbf{Model}: Soft-voting ensemble (LR, RF, XGBoost, GBoost)
    \item \textbf{Pipeline}: scikit-learn Pipeline with ColumnTransformer, scaling, encoding
    \item \textbf{Interface}: Enterprise dashboard with 4-page web application
    \item \textbf{API}: RESTful JSON prediction endpoint with full validation
\end{itemize}

% ══════════════════════════════════════════════════════════════════
% 2. PROJECT STRUCTURE
% ══════════════════════════════════════════════════════════════════
\section{Project Structure}

\begin{lstlisting}[style=bash]
E:\Student Depression\
  app.py                                       # Flask application (261 lines)
  student_depression_ensemble_v3_optimized.pkl # Trained ML pipeline (62.5 MB)
  templates/
    index.html          # Assessment form (1504 lines, 49 KB)
    result.html         # Prediction result (459 lines, 19 KB)
    about.html          # Model architecture (292 lines, 15 KB)
    documentation.html  # API reference (360 lines, 20 KB)
  Student Depression Dataset - Cleaned.csv     # Training data (2196 KB)
  Student Depression.ipynb                     # Model training notebook
  *.png                                        # Training visualizations
  PROJECT_REPORT.md                            # Markdown report
  PROJECT_REPORT.tex                           # LaTeX report (this file)
\end{lstlisting}

% ══════════════════════════════════════════════════════════════════
% 3. MACHINE LEARNING PIPELINE
% ══════════════════════════════════════════════════════════════════
\section{Machine Learning Pipeline}

\subsection{Pipeline Architecture}

The model is serialized as a single \texttt{sklearn.pipeline.Pipeline} object containing a \texttt{ColumnTransformer} preprocessor followed by a soft-voting \texttt{VotingClassifier}.

\begin{lstlisting}[style=bash]
Pipeline
  preprocessor (ColumnTransformer)
    num (27 columns: 11 raw + 16 engineered)
      SimpleImputer(strategy='median')
      StandardScaler()
    ord (1 column: Sleep Duration)
      SimpleImputer(strategy='most_frequent')
      OrdinalEncoder(handle_unknown='use_encoded_value')
    nom (4 columns: City, Profession, Dietary Habits, Degree)
      SimpleImputer(strategy='most_frequent')
      OneHotEncoder(handle_unknown='ignore')
  classifier (VotingClassifier, soft voting)
    LogisticRegression (C=0.1, solver=lbfgs)              weight = 1.0
    RandomForestClassifier (n=400, max_depth=12)          weight = 2.0
    XGBClassifier (n=200, max_depth=4, lr=0.05)           weight = 2.0
    GradientBoostingClassifier (n=300, max_depth=5, lr=0.05) weight = 1.5
\end{lstlisting}

\subsection{Raw Input Features (16)}

\renewcommand{\arraystretch}{1.2}
\begin{longtable}{p{3.5cm}p{2.5cm}p{2.5cm}p{5cm}}
\hline
\textbf{Name} & \textbf{Type} & \textbf{Range} & \textbf{Input Method} \\
\hline
Gender & Binary & 0 = Female, 1 = Male & Dropdown \\
Age & Integer & 10--50 & Number input \\
City & Text & Free entry & Text input \\
Profession & Hidden & Always ``Student'' & Hidden input \\
Academic Pressure & Integer & 1--5 & Radio group \\
Work Pressure & Integer & 1--5 & Radio group \\
CGPA & Float & 0.00--4.00 & Number input \\
Study Satisfaction & Integer & 1--5 & Radio group \\
Job Satisfaction & Integer & 1--5 & Radio group \\
Sleep Duration & Categorical & 4 options & Dropdown \\
Dietary Habits & Categorical & 4 options & Dropdown \\
Degree & Text & Free entry & Text input \\
Suicidal Thoughts & Binary & 0 = No, 1 = Yes & Dropdown \\
Work/Study Hours & Integer & 0--20 & Number input \\
Financial Stress & Integer & 1--5 & Radio group \\
Family History of MI & Binary & 0 = No, 1 = Yes & Dropdown \\
\hline
\end{longtable}

\subsection{Engineered Features (16)}

\begin{longtable}{p{4.5cm}p{5.5cm}p{4.5cm}}
\hline
\textbf{Feature} & \textbf{Formula} & \textbf{Purpose} \\
\hline
Total\_Pressure & AP + WP & Cumulative stress load \\
Pressure\_Satisfaction\_Ratio & (AP + WP) / (SS + JS + 1) & Pressure relative to satisfaction \\
Sleep\_Deprived & 1 if ``Less than 5 hours'' & Binary deprivation flag \\
Sleep\_Quality & Ordinal map (4 levels) & Quality score (0--3) \\
High\_Risk & Suicidal $\land$ (FS $\geq$ 4) & High-risk interaction \\
Stress\_Multiplier & AP $\times$ WP & Compounding stress \\
Overall\_Satisfaction & SS + JS & Total satisfaction \\
Satisfaction\_Pressure\_Ratio & (SS + JS + 1) / (AP + WP + 1) & Inverse satisfaction-pressure \\
WorkLife\_Balance & (SS + JS) / (Hours + 1) & Balance metric \\
FinAcad\_Stress\_Interaction & FS $\times$ AP & Double stress interaction \\
Suicide\_Risk\_Severity & Suicidal $\times$ (FS + AP + WP) / 2 & Weighted risk severity \\
MH\_Vulnerability & Weighted composite (5 factors) & Multi-factor vulnerability \\
CGPA\_Performance\_Gap & (4.0 -- CGPA) $\times$ AP & Academic distress \\
Perfect\_Health & All favorable conditions & Absence of risk markers \\
Extreme\_Stress & $\geq$ 3 of 5 extreme conditions & Multi-domain extreme stress \\
CGPA\_Achievement & CGPA / (AP + 1) & Performance under pressure \\
\hline
\end{longtable}

\noindent\textit{Abbreviations:} AP = Academic Pressure, WP = Work Pressure, SS = Study Satisfaction, JS = Job Satisfaction, FS = Financial Stress.

\subsection{Performance}

\begin{itemize}[nosep]
    \item \textbf{AUC-ROC}: 99.2\%
    \item \textbf{Training samples}: ~27,900
    \item \textbf{Cross-validation}: Stratified K-fold
    \item \textbf{Voting strategy}: Soft (probability-weighted)
    \item \textbf{Class weights}: [1.0, 2.0, 2.0, 1.5] across the four models
\end{itemize}

% ══════════════════════════════════════════════════════════════════
% 4. WEB APPLICATION
% ══════════════════════════════════════════════════════════════════
\section{Web Application Architecture}

\subsection{Technology Stack}

\begin{tabularx}{\textwidth}{l X}
\hline
\textbf{Component} & \textbf{Technology} \\
\hline
Backend & Python 3.13, Flask 3.x \\
Frontend & HTML5, CSS3 (1500+ lines), JavaScript (vanilla) \\
Model Runtime & scikit-learn 1.x, XGBoost \\
Serialization & joblib \\
\hline
\end{tabularx}

\subsection{Routes}

\begin{tabularx}{\textwidth}{l l l X}
\hline
\textbf{Route} & \textbf{Method} & \textbf{Description} & \textbf{Response} \\
\hline
\texttt{/} & GET & Assessment form & \texttt{index.html} \\
\texttt{/predict} & POST & Form submit $\rightarrow$ prediction & \texttt{result.html} \\
\texttt{/about} & GET & Model architecture \& metrics & \texttt{about.html} \\
\texttt{/documentation} & GET & API reference \& usage & \texttt{documentation.html} \\
\texttt{/api/predict} & POST & JSON API $\rightarrow$ prediction & JSON \\
\hline
\end{tabularx}

\subsection{Request Flow}

\begin{enumerate}[nosep]
    \item User loads \texttt{/} $\rightarrow$ \texttt{index.html} rendered with assessment form
    \item User fills 15 visible fields (progress bar tracks completion)
    \item Client-side validation runs in real-time (age, CGPA, hours)
    \item On submit, loading overlay displays while server processes
    \item Server validates all 16 fields for presence, type, and range
    \item \texttt{engineer\_features()} derives 16 additional columns
    \item \texttt{model.predict()} and \texttt{model.predict\_proba()} execute
    \item \texttt{result.html} renders with prediction, probability, risk level, advice
    \item JSON API (\texttt{/api/predict}) follows the same pipeline but returns JSON
\end{enumerate}

\subsection{Validation}

\begin{itemize}[nosep]
    \item \textbf{Client-side}: Real-time validation on Age (10--50), CGPA (0.00--4.00), Work/Study Hours (0--20) with success/error indicators under each field
    \item \textbf{Server-side}: All 16 fields checked for presence, non-empty values, numeric type correctness, and range bounds
    \item \textbf{API}: Returns HTTP 400 with descriptive error messages for missing fields, empty values, type errors, and out-of-range inputs
\end{itemize}

\subsection{Risk Classification}

\begin{tabularx}{\textwidth}{l X}
\hline
\textbf{Depression Probability} & \textbf{Risk Level} \\
\hline
$< 25\%$ & VERY LOW RISK \\
$25\%$ -- $49\%$ & LOW RISK \\
$50\%$ -- $74\%$ & MODERATE RISK \\
$\geq 75\%$ & HIGH RISK \\
\hline
\end{tabularx}

% ══════════════════════════════════════════════════════════════════
% 5. USER INTERFACE
% ══════════════════════════════════════════════════════════════════
\section{User Interface}

\subsection{Design System}

\begin{itemize}[nosep]
    \item \textbf{Theme}: Light enterprise dashboard
    \item \textbf{Colors}: Indigo primary (\#6366f1), slate text (\#1e293b), semantic success/danger/warning
    \item \textbf{Typography}: Inter (Google Fonts), weights 300--700
    \item \textbf{Layout}: Fixed sidebar (260px) + fluid main content
    \item \textbf{Responsiveness}: 4 breakpoints (desktop / 1024px / 768px / 480px)
    \item \textbf{Cross-browser}: \texttt{:has()} CSS with JavaScript \texttt{.checked} class fallback
\end{itemize}

\subsection{Form Pages}

\textbf{index.html} --- 3-section assessment form:

\begin{enumerate}[nosep]
    \item \textbf{Personal Information}: Gender (select), Age (number), City (text), Profession (hidden = ``Student'')
    \item \textbf{Academic \& Work}: Degree (text), CGPA (number), 5 radio groups (1--5) for Academic Pressure, Study Satisfaction, Work Pressure, Job Satisfaction, Work/Study Hours (number)
    \item \textbf{Lifestyle \& Health}: Sleep Duration (select), Dietary Habits (select), Financial Stress radio group, Suicidal Ideation (select), Family History (select)
\end{enumerate}

Key UI features:
\begin{itemize}[nosep]
    \item Progress bar tracking 15 visible fields (radio groups counted as one)
    \item Model banner showing ensemble architecture and 99\% AUC-ROC
    \item Google Forms-style radio button groups with circle indicators and labels
    \item Real-time validation with green checkmarks and red error indicators
    \item Loading overlay with animated progress bar during submission
    \item localStorage auto-save (disabled by default, ready to enable)
    \item Mobile sidebar toggle via fixed-position hamburger button
\end{itemize}

\subsection{Result Page}

\textbf{result.html} displays:
\begin{itemize}[nosep]
    \item Primary outcome: ``Depression detected'' or ``No depression detected''
    \item Color-coded risk level badge (red for HIGH, orange for MODERATE, green for LOW/VERY LOW)
    \item Animated probability bar with percentage label
    \item Metadata grid showing model confidence and risk classification
    \item Contextual advice list (clinical recommendations or preventive guidance)
    \item Action buttons: ``New Assessment'' and ``Print Report''
\end{itemize}

\subsection{Information Pages}

\begin{itemize}[nosep]
    \item \textbf{about.html}: Ensemble architecture (4 models), hyperparameters, performance metrics (99.2\% AUC-ROC), feature engineering descriptions, limitations disclaimer
    \item \textbf{documentation.html}: Web interface guide, REST API reference with full JSON schema and curl example, input field table with types/ranges, risk classification table, validation rules
\end{itemize}

\subsection{Responsive Breakpoints}

\begin{tabularx}{\textwidth}{l X}
\hline
\textbf{Screen Width} & \textbf{Behavior} \\
\hline
$>$ 1024px & 2-column form grid, sidebar pinned left \\
768--1024px & Reduced padding, model-divider hidden \\
$<$ 768px & Sidebar slides off-screen, hamburger toggle, 1-column form, radio options shrink \\
$<$ 480px & Compact model banner, tighter radio groups (48px min), smaller text \\
\hline
\end{tabularx}

% ══════════════════════════════════════════════════════════════════
% 6. API REFERENCE
% ══════════════════════════════════════════════════════════════════
\section{API Reference}

\subsection{POST /api/predict}

\textbf{Request} (JSON):

\begin{lstlisting}[style=py]
{
  "Gender": 1,
  "Age": 22,
  "City": "Karachi",
  "Profession": "Student",
  "Academic Pressure": 3,
  "Work Pressure": 2,
  "CGPA": 3.2,
  "Study Satisfaction": 4,
  "Job Satisfaction": 3,
  "Sleep Duration": "7-8 hours",
  "Dietary Habits": "Moderate",
  "Degree": "BSc Computer Science",
  "Have you ever had suicidal thoughts ?": 0,
  "Work/Study Hours": 8,
  "Financial Stress": 2,
  "Family History of Mental Illness": 0
}
\end{lstlisting}

\textbf{Response} (200):

\begin{lstlisting}[style=py]
{
  "prediction": 0,
  "depression_probability": 0.1308,
  "no_depression_probability": 0.8692
}
\end{lstlisting}

\textbf{Error} (400 / 500):

\begin{lstlisting}[style=py]
{ "error": "descriptive message" }
\end{lstlisting}

\textbf{curl example}:

\begin{lstlisting}[style=bash]
curl -X POST http://localhost:5000/api/predict \
  -H "Content-Type: application/json" \
  -d '{"Gender":1,"Age":22,"City":"Karachi","Profession":"Student",
       "Academic Pressure":3,"Work Pressure":2,"CGPA":3.2,
       "Study Satisfaction":4,"Job Satisfaction":3,
       "Sleep Duration":"7-8 hours","Dietary Habits":"Moderate",
       "Degree":"BSc CS","Have you ever had suicidal thoughts ?":0,
       "Work/Study Hours":8,"Financial Stress":2,
       "Family History of Mental Illness":0}'
\end{lstlisting}

% ══════════════════════════════════════════════════════════════════
% 7. DEPLOYMENT
% ══════════════════════════════════════════════════════════════════
\section{Deployment}

\subsection{Local Setup}

\begin{lstlisting}[style=bash]
pip install flask pandas numpy scikit-learn xgboost joblib
python app.py
# Application starts at http://localhost:5000
\end{lstlisting}

\subsection{Environment Variables}

\begin{tabularx}{\textwidth}{l l X}
\hline
\textbf{Variable} & \textbf{Default} & \textbf{Description} \\
\hline
\texttt{PORT} & 5000 & Server port number \\
\hline
\end{tabularx}

\subsection{Dependencies}

\begin{itemize}[nosep]
    \item flask (3.x)
    \item pandas (2.x)
    \item numpy (1.x)
    \item scikit-learn (1.x)
    \item xgboost (2.x)
    \item joblib (1.x)
\end{itemize}

% ══════════════════════════════════════════════════════════════════
% 8. DEVELOPMENT NOTES
% ══════════════════════════════════════════════════════════════════
\section{Development Notes}

\begin{itemize}[nosep]
    \item \texttt{OneHotEncoder} uses \texttt{handle\_unknown='ignore'} --- unseen City/Degree/Profession values produce a zero vector rather than errors, making free-text inputs safe.
    \item \texttt{OrdinalEncoder} for Sleep Duration uses \texttt{handle\_unknown='use\_encoded\_value'} for robustness against unexpected categorical values.
    \item Feature names with spaces (e.g., \texttt{"Have you ever had suicidal thoughts ?"}) are preserved exactly to match the training data column names.
    \item The \texttt{submit-bar} at the form bottom displays ``32 clinical features including 16 engineered indicators'' to communicate model sophistication to users.
    \item All templates include a \texttt{<meta name="viewport">} tag and a consistent set of CSS custom properties (\texttt{--primary}, \texttt{--text}, \texttt{--border}, etc.) for theme maintainability.
    \item The \texttt{:has()} CSS selector for radio button styling is paired with a JavaScript \texttt{.checked} class toggle for compatibility with browsers lacking native \texttt{:has()} support (Firefox $<$ 121).
    \item Radio button progress tracking counts each group as a single field (using a JavaScript \texttt{Set} keyed by \texttt{input[name]}), so the 5 rating groups contribute 5 toward the 15-field total rather than 25.
\end{itemize}

% ══════════════════════════════════════════════════════════════════
% END
% ══════════════════════════════════════════════════════════════════
\end{document}
